\subsection{LINEAR EQUATIONS}

Let E, F be two A-modules. Every equation of the form \(u(x) = y_{0}\) , where \(u: E \to F\) is a given linear mapping, \(y_{0}\) a given element of F and the unknown x is subjected to the condition that it takes its values in E, is called a linear equation; \(y_{0}\) is called the right hand side of the equation; if \(y_{0} = 0\) , the equation is called homogeneous linear.

Every element \(x_{0} \in E\) such that \(u(x_{0}) = y_{0}\) is called a solution of the linear equation \(u(x) = y_{0}\) .†

It is often said, loosely speaking, that a problem is linear if it is equivalent to determining the solutions of a linear equation.

Given a linear equation \(u(x) = y_{0}\) , the equation \(u(x) = 0\) is called the homogeneous linear equation associated with \(u(x) = y_{0}\) .

PROPOSITION 14. If \(x_{0}\) is a solution of the linear equation \(u(x) = y_{0}\) , the set of solutions of this equation is equal to the set of elements \(x_{0} + z\) , where z runs through the set of solutions of the associated homogeneous equation \(u(x) = 0\) .

The relation \(u(x) = y_0\) may be written as \(u(x) = u(x_0)\) , which is equivalent to \(u(x - x_0) = 0\) .

In other words, if the equation \(u(x) = y_{0}\) has at least one solution \(x_{0}\) , the set of its solutions is the set \(x_{0} + u^{-1}(0)\) , obtained by translation from the kernel \(u^{-1}(0)\) of u. Observe that \(u^{-1}(0)\) , being a submodule, is never empty, since it contains 0 (called the zero solution, or trivial solution, of the homogeneous equation \(u(x) = 0\) ).

By virtue of Proposition 14, for a linear equation \(u(x) = y_0\) to have exactly one solution, it is necessary and sufficient that it have at least one solution and that \(u^{-1}(0) = \{0\}\) (in other words, that the associated homogeneous equation have no non-zero solution, or also that \(u\) be injective); in this case, for all \(y \in F\) , the equation \(u(x) = y\) has at most one solution.

PROPOSITION 15. Let \(u\) be a linear mapping of a module \(\mathbf{E}\) into a module \(\mathbf{F}\) . If the equation \(u(x) = y_0\) has at least one solution, \(y_0\) is orthogonal to the kernel of \(^t u\) .

To say that \(u(x) = y_{0}\) admits a solution means that \(y_{0} \in u(\mathbf{E})\) and the proposition follows from no. 5, Corollary to Proposition 8.

Observe that the necessary criterion for the existence of a solution of \(u(x) = y_{0}\) , given by Proposition 15, is sufficient when A is a field (§ 7, no. 6, Proposition 12), but not in general (Exercise 10).

Remarks. (1) Let E be an A-module, \((\mathbf{F}_{i})_{i\in I}\) a family of A-modules and for all \(i\in I\) let \(u_{i}:E\to F_{i}\) be a linear mapping. Every system of linear equations

\[
u _ {i} (x) = y _ {i} \quad (i \in I)\tag{29}
\]

where the \(y_{i} \in \mathbf{F}_{i}\) are given, is equivalent to a single linear equation \(u(x) = y\) , where \(u\) is the mapping \(x \mapsto (u_{i}(x))\) of \(\mathbf{E}\) into \(\mathbf{F} = \prod_{i \in I} \mathbf{F}_{i}\) and \(y = (y_{i})\) . The system (29) is called homogeneous if \(y_{i} = 0\) for all \(i \in I\) .

\begin{enumerate}
\def\labelenumi{(\arabic{enumi})}
\setcounter{enumi}{1}
\tightlist
\item
  Suppose that E admits a basis \((a_{\lambda})_{\lambda \in L}\) ; if we set \(u(a_{\lambda}) = b_{\lambda}\) for all \(\lambda \in L\) , to say that \(x = \sum_{\lambda \in L} \xi_{\lambda} a_{\lambda}\) satisfies the equation \(u(x) = y_0\) is equivalent to saying that the family (of finite support) \((\xi_{\lambda})_{\lambda \in L}\) of elements of A satisfies the relation
\end{enumerate}

\[
\sum_ {\lambda \in L} \xi_ {\lambda} b _ {\lambda} = y _ {0}.\tag{30}
\]

Conversely, looking for families \((\xi_{\lambda})_{\lambda \in L}\) of elements of A of finite support satisfying (30), is equivalent to solving the linear equation \(u(x) = y_0\) , where \(u\) is the unique linear mapping of E into F such that \(u(a_{\lambda}) = b_{\lambda}\) for all \(\lambda \in L\) (§ 1, no. 11, Corollary 3 to Proposition 17).

\begin{enumerate}
\def\labelenumi{(\arabic{enumi})}
\setcounter{enumi}{2}
\tightlist
\item
  A linear equation \(u(x) = y_0\) is called scalar when \(F = A_s\) and therefore \(u\) is a linear form on \(E\) and \(y_0\) a scalar. If \(E\) admits a basis \((a_\lambda)_{\lambda \in L}\) , it follows from Remark (2) that such an equation may also be written as
\end{enumerate}

\[
\sum_ {\lambda \in L} \xi_ {\lambda} \alpha_ {\lambda} = y _ {0} \in A\tag{31}
\]

where the family of scalars \((\alpha_{\lambda})\) is arbitrary and where it is understood that the family \((\xi_{\lambda})\) must have finite support. In general, by the solution (in A) of a system of scalar linear equations

\[
\sum_ {\lambda \in L} \xi_ {\lambda} \alpha_ {\lambda i} = \eta_ {i} \quad (i \in I)\tag{32}
\]

where \(\alpha_{\lambda i} \in A\) and \(\eta_i \in A\) , is understood a family \((\xi_\lambda)_{\lambda \in L}\) of elements of \(A\) of finite support and satisfying (32); the \(\alpha_{\lambda i}\) are called the coefficients of the system of equations and the \(\eta_i\) the right hand sides. The solution of such a system is equivalent to that of the equation \(u(x) = y\) , where \(y = (\eta_i)\) and \(u: A_s^{(L)} \to A_s^I\) is the linear mapping

\[
\left(\xi_ {\lambda}\right) \mapsto \left(\sum_ {\lambda \in L} \xi_ {\lambda} \alpha_ {\lambda i}\right).
\]

\begin{enumerate}
\def\labelenumi{(\arabic{enumi})}
\setcounter{enumi}{3}
\tightlist
\item
  A linear system (32) is also called a system of left scalar linear equations when it is necessary to avoid confusion. A system of equations
\end{enumerate}

\[
\sum_ {\lambda \in L} \alpha_ {\lambda i} \xi_ {\lambda} = \eta_ {i} \quad (i \in I)\tag{33}
\]

is likewise called a system of right scalar linear equations; such a system can immediately be transformed into a system (32) by considering the \(\xi_{\lambda}, \eta_{i}\) and \(\alpha_{\lambda i}\) as belonging to the opposite ring A \(^{0}\) to A.

